\chapter{Coronary Angioplasty}

\section*{Overview}
Coronary angioplasty is a minimally invasive procedure used to open clogged heart arteries. It utilizes a balloon-tipped catheter to widen the lumen and place a drug-eluting stent. The exact approach, setting, and preparation depend on the indication, anatomy, and the patient's health.

\section*{Alternative Names}
Percutaneous Coronary Intervention, PCI, Balloon angioplasty

\section*{Principle}
A balloon-tipped catheter is positioned across a narrowed coronary artery and inflated to crush the plaque and stretch the vessel wall. A stent is usually deployed to hold the artery open and prevent elastic recoil and restenosis.
\section*{History}
Andreas Gruentzig performed the first successful balloon coronary angioplasty in Zurich in 1977. Coronary stents were introduced in the late 1980s.

\section*{Why It Is Done}
Ordered for patients experiencing acute myocardial infarction (STEMI/NSTEMI), unstable angina, or stable angina refractory to medical therapy.

\section*{Is This a Primary or Secondary Test or Procedure}
Primary emergency treatment for acute ST-elevation myocardial infarction (STEMI).

\section*{How It Is Performed}
Under local anesthesia and conscious sedation. A catheter is inserted through the radial or femoral artery and guided to the heart. A balloon is inflated at the blockage site, and a stent is deployed.

\section*{Preparation}
Fasting, IV access, and baseline testing are arranged according to urgency and the planned sedation. The cardiology team selects antiplatelet and anticoagulant treatment based on the indication, whether a stent is used, kidney function, and bleeding risk; take or stop medicines only according to that plan.

\section*{Contraindications and Precautions}
None in acute life-threatening MI. Relative contraindications include severe renal failure (due to contrast) or active bleeding.

\section*{Risks}
Coronary artery dissection or rupture, acute MI, contrast-induced nephropathy, bleeding at the access site, or stent thrombosis.

\section*{Aftercare and Recovery}
Monitor vital signs and the puncture site for bleeding. Radially accessed patients may be able to sit up sooner, according to the team's instructions. Take prescribed antiplatelet or anticoagulant treatment exactly as directed and do not stop it without contacting cardiology.

\section*{Outcome and Follow-up}
Successful angioplasty restores normal coronary blood flow (grade 3 TIMI flow) and resolves chest pain/ECG changes.

\section*{Expected Results}
TIMI Grade 3 flow in the target artery; stent fully expanded; stenosis reduced to <10%.

\section*{Follow-up}
Cardiology follow-up, medication review, risk-factor treatment, and cardiac rehabilitation may be recommended. The duration and combination of antiplatelet or anticoagulant treatment are individualized rather than fixed for every patient.

\section*{When to Seek Medical Care}
Follow the procedure team's specific instructions. Seek urgent medical assessment for severe or worsening pain, uncontrolled bleeding, fever or signs of infection, trouble breathing, chest pain, fainting, new weakness or confusion, inability to urinate, or any rapidly worsening symptom. The appropriate warning signs depend on the procedure and the patient's medical condition.

\section*{References}
\begin{enumerate}
  \item MedlinePlus. Coronary Angioplasty. U.S. National Library of Medicine; 2025.
  \item Current Medical Diagnosis and Treatment. McGraw-Hill; 2025.
\end{enumerate}

\clearpage
